% =============================================================================
% Appendix A: Software Repository and Pipeline
% =============================================================================
\mychapter{Appendix A --- Software Repository and Pipeline}

This appendix documents the software repository structure, installation instructions, and pipeline
execution steps for the automated XCT defect detection pipeline developed in this thesis. The
repository structure below reflects the codebase's current, actual layout as of August 2026, which
has been substantially restructured since the April 2026 mid-term submission (a \url{src/}
package and a \url{scripts/} entry-point folder were introduced, and preprocessing/label-generation
logic moved out of the originally-described \url{data/loader.py} and
\url{data/pseudo_labels.py}); see the Note on this combined document.

\mysection{A.1 Repository Access}
\begin{table}[htbp]
\centering
\caption{Repository information.}
\begin{tabular}{lp{10cm}}
\toprule
\textbf{Field} & \textbf{Value} \\
\midrule
Platform & GitHub \\
Repository URL & \url{https://github.com/jobgkj/podfam_research_project_XCT_Anlaysis} \\
Branch & main \\
Language & Python 3.10+ \\
Last commit & August 2026 \\
\bottomrule
\end{tabular}
\end{table}

\mysection{A.2 Repository Structure (current)}
Top level:
\begin{itemize}
\item \texttt{config.py} --- all hyperparameters and paths, including PODFAM-phase additions
  (\param{TRAINING\_SAMPLES\_OVERRIDE}, \param{MAX\_TRAINING\_SAMPLES}, \param{EXPERIMENT\_NAME},
  Bernsen parameters)
\item \texttt{pipeline.py} --- end-to-end 2D(+3D) training entry point
\end{itemize}

\url{src/} --- preprocessing, thresholding, and I/O (introduced after the mid-term; supersedes
the originally-planned \url{data/loader.py} and \url{data/pseudo_labels.py}):
\begin{itemize}
\item \url{src/preprocess.py} --- per-slice preprocessing (median filter + adaptive NLM
  denoising)
\item \url{src/thresholding.py} --- Otsu, Yen, and Bernsen thresholding (Bernsen is the
  PODFAM-phase primary pseudo-label source; see Part II, Section~9.3)
\item \url{src/sample_mask.py} --- circular sample-boundary detection
\item \url{src/io.py} --- mask generation orchestration
\item \url{src/metrics.py} --- per-pore connected-component statistics
\end{itemize}

\url{data/} --- dataset and caching:
\begin{itemize}
\item \url{data/cache.py} --- memory-mapped volume/mask cache (introduced after the mid-term, for
  training on volumes too large to hold fully in RAM)
\item \url{data/dataset.py}, \url{data/dataset_3d.py} --- PyTorch Dataset and patch
  extraction
\item \url{data/augmentation.py} --- online augmentation transforms
\end{itemize}

\url{models/}:
\begin{itemize}
\item \url{models/unet2d.py} --- 2D U-Net architecture
\item \url{models/unet3d.py} --- 3D U-Net architecture
\item \url{models/thesis_analysis.py} --- figure-generation script (renamed and moved from the
  originally-described root-level \texttt{thesis\_analysis.py})
\end{itemize}

\url{training/}:
\begin{itemize}
\item \url{training/losses.py} --- BCE, Dice, Focal, Dice-Focal
\item \url{training/metrics.py} --- Dice, IoU, Precision, Recall
\item \url{training/trainer.py} --- training loop with MLflow tracking, early stopping, and
  best-model checkpointing
\end{itemize}

\url{scripts/} --- command-line entry points (introduced after the mid-term; supersedes the
originally-described root-level \texttt{run\_preprocess.py}):
\begin{itemize}
\item \url{scripts/run_preprocess.py}, \url{scripts/run_single_sample.py},
  \url{scripts/generate_training_data.py} --- preprocessing and mask-generation drivers
\item \url{scripts/predict_and_visualize_3d.py} --- trained-model inference and 3D
  reconstruction
\item \url{scripts/compare_methods_report.py} --- classical-vs-U-Net comparison table generator
\item \url{scripts/check_training_status.py} --- MLflow-based training status checker
\end{itemize}

\url{data/raw/}, \url{data/processed/}, \url{data/masks/} --- input and intermediate data
(not in Git); \url{mlruns/} --- MLflow tracking store (not in Git); \url{artifacts/} --- trained
model checkpoints (not in Git); \url{results/} --- figures, predictions, comparison tables, and
the thesis chapter document.

\mysection{A.3 Installation}
\mysubsection{A.3.1 Prerequisites}
\begin{itemize}
\item Python 3.10 or later
\item Git
\item pip package manager
\item NVIDIA GPU with CUDA support (recommended for model training)
\end{itemize}

\mysubsection{A.3.2 Setup Steps}
Clone the repository and install dependencies:
\begin{verbatim}
git clone https://github.com/jobgkj/podfam_research_project_XCT_Anlaysis
cd podfam_research_project_XCT_Anlaysis/xct_defect_detection
python -m venv .venv
.venv\Scripts\activate      (Windows)
source .venv/bin/activate   (Linux/macOS)
pip install -r requirements.txt
\end{verbatim}

\mysubsection{A.3.3 Required Python Packages}
\begin{table}[htbp]
\centering
\caption{Key Python dependencies.}
\begin{tabular}{p{3.2cm}p{4.4cm}p{5.6cm}}
\toprule
\textbf{Package} & \textbf{Version} & \textbf{Purpose} \\
\midrule
torch & 2.0+ (cu124 build for GPU) & Deep learning framework \\
numpy & 1.24+ & Array operations \\
tifffile & 2023+ & TIFF stack loading \\
scikit-image & 0.21+ & Image processing, NLM denoising, marching cubes \\
scipy & 1.11+ & Median filter, polynomial fitting \\
matplotlib & 3.7+ & Figure generation \\
mlflow & 2.5+ & Experiment tracking \\
python-docx & 1.2+ & Thesis document generation \\
\bottomrule
\end{tabular}
\end{table}

\mysection{A.4 Running the Pipeline}
Preprocess and generate masks for one sample:
\begin{verbatim}
python scripts/run_single_sample.py sample_07
\end{verbatim}

Run the full 2D U-Net training pipeline:
\begin{verbatim}
python pipeline.py
\end{verbatim}

Run inference and 3D visualisation on a trained checkpoint:
\begin{verbatim}
python scripts/predict_and_visualize_3d.py sample_07
\end{verbatim}

Generate the classical-vs-U-Net comparison table:
\begin{verbatim}
python scripts/compare_methods_report.py
\end{verbatim}

View experiment results in MLflow:
\begin{verbatim}
mlflow ui --backend-store-uri mlruns
\end{verbatim}
then open \url{http://localhost:5000} in a browser.

\mysection{A.5 Version Control and Reproducibility}
\begin{itemize}
\item All source code, configuration files, and documentation are committed to the main branch
\item Raw data, preprocessed outputs, model checkpoints, and MLflow logs are excluded via
  \texttt{.gitignore} due to file size
\item All experiments are logged in MLflow with full hyperparameter and metric records, enabling
  exact reproduction of any training run, tagged by \texttt{EXPERIMENT\_NAME} (Section~9.4)
\end{itemize}

\mysection{A.6 Interactive 3D Defect Viewer (gui\_app.py)}
A standalone desktop GUI, \texttt{gui\_app.py}, provides an interactive front end to the trained U-Net
detector, wrapping the same inference pipeline as \url{scripts/predict_and_visualize_3d.py} in a
Tkinter application:

\begin{itemize}
\item \textbf{Folder selection}: a file dialog lets the user pick any local folder of TIFF slices
  (not limited to the pre-registered \texttt{sample\_NN} folders under \texttt{data/raw/}), and a
  second dialog selects the checkpoint to run, defaulting to
  \texttt{artifacts/best\_model\_\{EXPERIMENT\_NAME\}.pt}.
\item \textbf{Inference}: on ``Run'', the selected volume is preprocessed with the same
  \texttt{src/preprocess.py} pipeline used throughout this thesis, then run through the network using
  the identical sliding-window inference procedure and patch configuration ($256 \times 256$ patches,
  stride 128) as \texttt{scripts/predict\_and\_visualize\_3d.py}. Preprocessing and inference progress
  are shown live via a progress bar, running on a background thread so the window stays responsive.
\item \textbf{Defect-zone highlighting}: the predicted defect voxels are reconstructed into a
  marching-cubes mesh and rendered in red against a dark background; only defect voxels are drawn, so
  the mesh itself directly highlights the predicted defect zones in three dimensions, with no separate
  overlay step required.
\item \textbf{Interactive 3D controls}: the mesh is displayed in an embedded, mouse-controllable
  viewer, drag to rotate, scroll to zoom, plus a navigation toolbar for pan, zoom, reset, and
  saving the current view as an image.
\item \textbf{Metrics panel}: after each run, a metrics panel reports the total defect fraction, pore
  count, mean pore area, and mean equivalent diameter for the predicted volume, computed via
  \texttt{src/metrics.py}'s connected-component analysis pooled across every slice, the same
  computation used to produce Tables 10.1--10.5.
\end{itemize}

Run it from the repository root with:
\begin{verbatim}
python gui_app.py
\end{verbatim}
